\documentclass[12pt,letterpaper]{article}

%% ─── PACKAGES ────────────────────────────────────────────────────────────────
\usepackage[margin=1in]{geometry}
\usepackage{amsmath, amssymb, amsthm}
\usepackage{mathtools}
\usepackage{hyperref}
\usepackage{listings}
\usepackage{xcolor}
\usepackage{booktabs}
\usepackage{longtable}
\usepackage{enumitem}
\usepackage{fancyhdr}
\usepackage{titlesec}
\usepackage{abstract}
\usepackage{microtype}
\usepackage{datetime2}
\usepackage{cleveref}

%% ─── METADATA ────────────────────────────────────────────────────────────────
\hypersetup{
  pdftitle={Meta-Relativity and PIRTM: Defensive Publication and Comprehensive Technical Report},
  pdfauthor={Citizen Gardens \\ The Foundation of Multiplicity},
  pdfsubject={Prior Art Defensive Publication},
  pdfkeywords={Meta-Relativity, PIRTM, ZM operator stack, prime-indexed recursion,
               contractivity certification, Multiplicity Theory, AL-GFT,
               spin foam microfoundations, operator semigroups},
  colorlinks=true,
  linkcolor=black,
  citecolor=blue,
  urlcolor=blue
}

%% ─── THEOREM ENVIRONMENTS ────────────────────────────────────────────────────
\newtheorem{theorem}{Theorem}[section]
\newtheorem{lemma}[theorem]{Lemma}
\newtheorem{corollary}[theorem]{Corollary}
\newtheorem{proposition}[theorem]{Proposition}
\theoremstyle{definition}
\newtheorem{definition}[theorem]{Definition}
\newtheorem{axiom}[theorem]{Axiom}
\theoremstyle{remark}
\newtheorem{remark}[theorem]{Remark}


%% ── operators ─────────────────────────────────────────────
\DeclareMathOperator{\GapLB}{GapLB}
\DeclareMathOperator{\SlopeUB}{SlopeUB}
\DeclareMathOperator{\spec}{spec}
\DeclareMathOperator{\essran}{ess\,ran}
\DeclareMathOperator{\tr}{tr}
\DeclareMathOperator{\sgap}{\delta_S}
\newcommand{\PP}{\mathbb{P}}
\newcommand{\RR}{\mathbb{R}}
\newcommand{\CC}{\mathbb{C}}
\newcommand{\ZZ}{\mathbb{Z}}
\newcommand{\NN}{\mathbb{N}}
\newcommand{\HH}{\mathcal{H}}
\newcommand{\HS}{\mathcal{HS}}
\newcommand{\norm}[1]{\left\|#1\right\|}
\newcommand{\ip}[2]{\left\langle #1,\, #2 \right\rangle}
\newcommand{\abs}[1]{\left|#1\right|}
\newcommand{\lp}{\ell^2(\PP)}
\newcommand{\Lt}{L^2(\RR)}
\newcommand{\Cd}{\CC^d}
%% ─── CODE LISTING STYLE ──────────────────────────────────────────────────────
\lstset{
  language=Python,
  basicstyle=\ttfamily\footnotesize,
  keywordstyle=\color{blue}\bfseries,
  commentstyle=\color{gray}\itshape,
  stringstyle=\color{teal},
  numberstyle=\tiny\color{gray},
  numbers=left,
  frame=single,
  breaklines=true,
  captionpos=b,
  showstringspaces=false,
  tabsize=4
}

%% ─── HEADER / FOOTER ─────────────────────────────────────────────────────────
\pagestyle{fancy}
\fancyhf{}
\lhead{\footnotesize Meta-Relativity \& PIRTM --- Defensive Publication}
\rhead{\footnotesize Citizen Gardens, April 2026}
\cfoot{\thepage}

%% ─── TITLE ───────────────────────────────────────────────────────────────────
\title{%
  \textbf{Meta-Relativity and the Prime Interval Recursion Theory Machine (PIRTM):}\\[4pt]
  \large Defensive Prior Art Publication,\\
  Comprehensive Mathematical Overview, and\\
  Implementation Architecture Report\\[12pt]
  \normalsize \textit{Version 1.0 --- Defensive Publication}
}

\author{%
  Citizen Gardens \\ The Foundation of Multiplicity
}

\date{April 25, 2026}

%% ═══════════════════════════════════════════════════════════════════════════
\begin{document}
\maketitle
\thispagestyle{fancy}

%% ─── DEFENSIVE PUBLICATION NOTICE ──────────────────────────────────────────
\vspace{1em}
\noindent\fbox{\parbox{\textwidth}{%
  \textbf{Defensive Prior Art Publication Notice.}
  This document is published publicly by the Multiplicity Foundation for the
  express purpose of establishing prior art under 35 U.S.C.~\S\,102 and
  equivalent international provisions (EPC Art.~54). All mathematical
  frameworks, operator constructions, certification protocols, module
  architectures, and software designs described herein are disclosed as of
  the date of publication. This disclosure is intended to prevent any third
  party from obtaining patent claims over the described subject matter.
  The disclosed work is released under the \textbf{Prime Materia License}
  and Apache 2.0 where applicable. Nothing in this document constitutes a
  waiver of copyright.
}}
\vspace{1em}

\begin{abstract}
We present a comprehensive technical disclosure of the \emph{Meta-Relativity}
framework and its primary computational instantiation, the
\emph{Prime Interval Recursion Theory Machine} (PIRTM), developed under the
Multiplicity Foundation. Meta-Relativity is a prime-indexed, recursively
stable mathematical framework wherein physical and abstract structures are
modeled as prime-labeled operator semigroups acting on tensor-product
Hilbert spaces. The framework unifies operator-theoretic contractivity
certification, Zetafunction-inflected spectral analysis, and a compositional
operator architecture (the ZM universal operator
$U = A \otimes I \otimes I + I \otimes C \otimes I + I \otimes I \otimes E$
on $H = \ell^2(\mathbb{P}) \otimes L^2(\mathbb{R}) \otimes \mathbb{C}^d$)
with a software module stack whose correctness is enforced through ADR-governed
stub modules and formally-specified certification protocols.

This publication establishes priority for the following novel contributions:
(1) the ZM universal operator construction and its finite-prime truncation;
(2) the GapLB/SlopeUB contractivity certification protocol (Theorem 6 of the
ZM spec); (3) the \texttt{PIRTMPolicy} Protocol and goal-budget enforcement
mechanism; (4) the \texttt{ZMOperatorStack} module binding architecture;
(5) the co-dependency resolution pattern for \texttt{certify.py} and
\texttt{pirtm/zm/}; and (6) the five-gate ordered commit sequence for
provably-correct PIRTM deployment.
\end{abstract}

\tableofcontents
\newpage

%% ═══════════════════════════════════════════════════════════════════════════
\section{Executive Summary}
\label{sec:exec}

The \emph{Meta-Relativity} project, hosted at
\url{https://github.com/MultiplicityFoundation/Meta-Relativity}, constitutes
an original research and engineering program at the intersection of analytic
number theory, operator spectral theory, quantum gravity (spin foam
microfoundations), and formal software verification. The primary work product
is the \emph{Prime Interval Recursion Theory Machine} (PIRTM), a formally
specified computational runtime whose correctness contracts are derived from
first principles in prime-indexed Hilbert space theory.

\subsection*{Key Contributions Established by This Publication}

\begin{enumerate}[leftmargin=*]

\item \textbf{The ZM Universal Operator.}
A tensor-product decomposition
$U = A \otimes I + I \otimes C \otimes I + I \otimes I \otimes E$
on a three-component Hilbert space whose prime block $A = D_\sigma + K$
encodes Dirichlet-series weighting and Hilbert-Schmidt off-diagonal structure.

\item \textbf{GapLB/SlopeUB Certification.}
A four-step certification protocol (verify HS condition $\to$ check
multiplier norm $\to$ compute GapLB/SlopeUB $\to$ enforce lawfulness) whose
contractivity bound takes the closed form
\[
  \operatorname{GapLB}(w) \;\geq\;
  \inf_\theta \!\left[\delta_S(\theta) - 2\sum_p |w_p|\, b_p \right].
\]

\item \textbf{The PIRTMPolicy Protocol.}
A typed Protocol class enforcing the lawfulness inequality
$\sum_k \varepsilon_k < (1 - \gamma_i)\cdot\|P_i\|^{-1}$
at runtime via \texttt{goal\_budget} checking inside \texttt{iterate()}.

\item \textbf{The ZMOperatorStack Binding Architecture.}
A stub-first, ADR-governed module design pattern in which
\texttt{pirtm/zm/operator\_stack.py} names all binding surfaces before
implementation, co-committed with \texttt{pirtm/certify.py} to resolve
import co-dependency in a single atomic commit.

\item \textbf{Spin Foam as Prime-Gated Hilbert Space Instantiation.}
The identification of spin foam area eigenvalues as a specific instantiation
of prime weights in $\ell^2(\mathbb{P})$, establishing the integration
pathway between \texttt{Spin\_Foam\_Microfoundations.md} and the ZM
operator stack.

\item \textbf{The Five-Gate Ordered Commit Sequence.}
A formal dependency ordering for PIRTM provable correctness deployment,
constituting a novel software engineering artifact.

\end{enumerate}

%% ═══════════════════════════════════════════════════════════════════════════
\section{Background and Framework}
\label{sec:background}

\subsection{Multiplicity Theory}

\begin{definition}[Multiplicity Space]
A \emph{multiplicity space} $(X, \mathcal{P}, \mu)$ is a set $X$ equipped
with a prime-indexed measure $\mu \colon \mathcal{P} \times X \to \mathbb{R}_{\geq 0}$
where $\mathcal{P} = \{2, 3, 5, 7, 11, \ldots\}$ denotes the rational primes,
such that the recurrence of any element $x \in X$ is governed by the
prime factorization of its multiplicity index.
\end{definition}

Multiplicity Theory reinterprets mathematical objects as recursively generated
patterns of prime-labeled interactions. Sets and modules become relationally
governed multiplicity spaces with identity and behavior preserved through
recursive feedback loops across scales. The prime-indexed structure is not
merely a labeling convenience; it is the \emph{generative substrate} from
which emergent structures arise.

\subsection{Meta-Relativity: Core Thesis}

Meta-Relativity extends the multiplicity framework to physical law. Rather
than treating spacetime geometry as fixed background structure, Meta-Relativity
posits that geometric and physical relationships are themselves manifestations
of prime-indexed operator interactions. The \emph{lawfulness} of a physical
configuration is certified by contractivity of the governing operator semigroup.

\begin{definition}[Lawful Configuration]
A configuration $\psi \in H$ is \emph{lawful} under operator family
$\{\Xi(t)\}_{t \geq 0}$ if there exists $\gamma \in (0,1)$ such that
\[
  \|\Xi(t)\psi - \Xi(t)\phi\| \;\leq\; \gamma^t \|\psi - \phi\|
  \quad \forall\, \psi, \phi \in H,\; \forall\, t \geq 0.
\]
\end{definition}

The contractivity parameter $\gamma$ is the \emph{lawfulness margin}, and its
certified lower bound is the central object of the PIRTM certification protocol.

%% ═══════════════════════════════════════════════════════════════════════════
\section{The ZM Universal Operator Construction}
\label{sec:zm}

\subsection{Hilbert Space Architecture}

The ZM framework (``Zeta-Multiplicity'') operates on a three-component
tensor product Hilbert space:
\[
  H \;=\; \ell^2(\mathbb{P}) \;\otimes\; L^2(\mathbb{R}) \;\otimes\; \mathbb{C}^d,
\]
where:
\begin{itemize}[leftmargin=*]
  \item $\ell^2(\mathbb{P})$ is the separable Hilbert space of square-summable
        sequences indexed by the rational primes, serving as the
        \emph{prime block} state space;
  \item $L^2(\mathbb{R})$ is the standard $L^2$ space over the reals, hosting
        the \emph{continuous multiplier} component; and
  \item $\mathbb{C}^d$ is a finite-dimensional ``edge'' or environmental
        component of dimension $d$.
\end{itemize}

\subsection{The Universal Operator}

\begin{definition}[ZM Universal Operator]
\label{def:zm-universal}
The \emph{ZM universal operator} is the densely defined self-adjoint operator
\[
  U \;=\; A \otimes I \otimes I \;+\; I \otimes C \otimes I
          \;+\; I \otimes I \otimes E
\]
where $A$, $C$, $E$ act on $\ell^2(\mathbb{P})$, $L^2(\mathbb{R})$,
$\mathbb{C}^d$ respectively.
\end{definition}

\subsubsection{The Prime Block \texorpdfstring{$A = D_\sigma + K$}{A}}

The prime block decomposes as:
\[
  A \;=\; D_\sigma + K,
\]
where $D_\sigma$ is the diagonal operator with entries $p^{-\sigma}$ for
$p \in \mathbb{P}$ (a Dirichlet-series weighting at real parameter
$\sigma > \frac{1}{2}$), and $K$ is the off-diagonal Hilbert-Schmidt
perturbation with kernel:
\[
  K(p, q) \;=\; p^{-\alpha}\, q^{-\alpha}\, h\!\left(\log p - \log q\right),
  \quad p,q \in \mathbb{P},
\]
for some $h \in L^2(\mathbb{R})$ with $\|h\|_{L^2} < \infty$, ensuring $K$
is Hilbert-Schmidt (HS).

\begin{proposition}[Hilbert-Schmidt Condition]
$K$ is Hilbert-Schmidt on $\ell^2(\mathbb{P})$ if and only if
\[
  \sum_{p,q \in \mathbb{P}} |K(p,q)|^2
  \;=\; \sum_{p,q} p^{-2\alpha} q^{-2\alpha} |h(\log p - \log q)|^2
  \;<\; \infty,
\]
which holds for $\alpha > \frac{1}{2}$ and $h \in L^2(\mathbb{R})$.
\end{proposition}

\subsubsection{The Multiplier Operator $C$}

The operator $C \colon L^2(\mathbb{R}) \to L^2(\mathbb{R})$ acts as a
Fourier multiplier:
\[
  C \;=\; \mathcal{F}^{-1} M_m \mathcal{F},
\]
where $M_m$ denotes multiplication by the real-valued symbol $m \in L^\infty(\mathbb{R})$
with essential range $[m_{\min}, m_{\max}]$. The multiplier range is the
parameter computed by \texttt{multiplier\_range()} in the ZM operator stack.

\subsubsection{The Recursive Operator \texorpdfstring{$\Xi(t)$}{Xi(t)}}

The full time-evolved operator family is:
\[
  \Xi(t) \;=\; \sum_{k=0}^{\infty} \alpha_k(t)\, T^k,
\]
where $T$ is a bounded contraction on $H$ and $\{\alpha_k(t)\}_{k \geq 0}$
are time-dependent coefficients satisfying $\sum_k |\alpha_k(t)| \leq 1$ for
all $t$. This is the bounded contraction named as \texttt{iterate()} in
\texttt{pirtm/core/recurrence.py}.

%% ═══════════════════════════════════════════════════════════════════════════
\section{Contractivity Certification Protocol}
\label{sec:cert}

\subsection{The Four-Step Protocol}

The ZM Certification Protocol (Section 6.2 of \texttt{ZM\_Meta\_Relativity.md})
consists of:

\begin{enumerate}[label=\textbf{Step \arabic*.}, leftmargin=*]
  \item \textbf{Verify HS condition.}
        Confirm $\sum_{p,q} |K(p,q)|^2 < \infty$
        (Hilbert-Schmidt norm finite).
  \item \textbf{Check multiplier norm.}
        Compute $[m_{\min}, m_{\max}]$ from the Fourier symbol of $C$.
  \item \textbf{Compute GapLB and SlopeUB.}
        Evaluate the contractivity margin bounds per Theorem 6.
  \item \textbf{Enforce lawfulness.}
        Assert $\operatorname{GapLB}(w) \geq \gamma_{\min}$.
\end{enumerate}

\subsection{Theorem 6: Contractivity Bounds}

\begin{theorem}[GapLB/SlopeUB Bounds]
\label{thm:gaplb}
Let $\{b_p\}_{p \in \mathbb{P}}$ be per-prime Lipschitz bounds,
$\{L_p\}_{p \in \mathbb{P}}$ be per-prime slope bounds, and
$\{w_p\}_{p \in \mathbb{P}}$ be the weight budget with
$w_{\mathrm{budget}} = \sum_p |w_p|$. Then:
\[
  \operatorname{GapLB}(w) \;\geq\;
  \inf_\theta \!\left[\delta_S(\theta) - 2\sum_p |w_p|\, b_p \right],
\]
\[
  \operatorname{SlopeUB}(w) \;=\; L \cdot w_{\mathrm{budget}},
\]
where $\delta_S(\theta) > 0$ is the spectral gap of $A$ at phase $\theta$,
$b_p = b_p(\alpha)$ is the Lipschitz bound derived from the HS kernel
parameter $\alpha$, and $L = \sup_p L_p < \infty$.
\end{theorem}

\begin{remark}
The mapping $\alpha \mapsto b_p$ is given by
$b_p \approx p^{-\alpha} \cdot \|h\|_{L^2}$
(ZM spec, \S2.3). This is the binding surface between the
\texttt{ZMOperatorStack} constructor parameter \texttt{alpha}
and the per-prime argument \texttt{b\_p} accepted by \texttt{certify()}.
Naming this mapping is a prerequisite for Day 2 implementation.
\end{remark}

\subsection{The Lawfulness Inequality}

\begin{definition}[PIRTMPolicy Lawfulness Constraint]
\label{def:lawfulness}
A \texttt{PIRTMPolicy} instance is \emph{lawful} if and only if
\[
  \sum_{k} \varepsilon_k
  \;<\;
  (1 - \gamma_i) \cdot \|P_i\|^{-1},
\]
where $\varepsilon_k$ are step-level perturbation magnitudes, $\gamma_i$ is
the contractivity parameter for module $i$, and $\|P_i\|$ is the operator
norm of the $i$-th projection. Violation of this constraint must raise a
\texttt{ValueError} inside \texttt{iterate()} at runtime.
\end{definition}

%% ═══════════════════════════════════════════════════════════════════════════
\section{PIRTM Software Architecture}
\label{sec:arch}

\subsection{Repository Structure}

The canonical implementation resides on the \texttt{Multiplicity} branch of
\texttt{MultiplicityFoundation/Meta-Relativity}. The top-level module
directory \texttt{pirtm/} contains the following active submodules:

\begin{longtable}{lll}
\toprule
\textbf{Module} & \textbf{Primary Abstraction} & \textbf{Status} \\
\midrule
\texttt{pirtm/sigma/} & Asymmetric kernel, ledger, MCP adapter & Implemented \\
\texttt{pirtm/spectral/} & Spectral operators, supermodule & Partially implemented \\
\texttt{pirtm/core/} & \texttt{recurrence.py}, \texttt{iterate()} & Partially implemented \\
\texttt{pirtm/bindings/} & Module binding surfaces & Partially implemented \\
\texttt{pirtm/channels/} & Channel protocols & Stub \\
\texttt{pirtm/backend/} & Backend adapters & Stub \\
\texttt{pirtm/mlir/} & MLIR dialect integration & Partial \\
\texttt{pirtm/zm/} & ZM universal operator stack & \textbf{Not yet created} \\
\texttt{pirtm/certify.py} & Certification protocol & \textbf{Not yet created} \\
\bottomrule
\caption{PIRTM module inventory as of April 2026.}
\end{longtable}

\subsection{Key Identified Gaps (Disclosure)}

The following structural gaps are disclosed as part of this prior art record:

\begin{enumerate}[leftmargin=*]
  \item \textbf{\texttt{certify\_state()} proxy inversion.}
        The current stub in \texttt{pirtm/core/} accepts operator
        arguments $(\Xi, \Lambda)$ but produces lawfulness certificates that
        are \emph{inverted} under transient conditions: the certificate passes
        for near-unstable operators and fails for stable ones. The correct
        signature is \texttt{certify\_state(b\_p, L\_p, w\_p, gamma\_min)}.

  \item \textbf{Missing \texttt{StepInfo.margin} field.}
        The \texttt{StepInfo} dataclass returned by \texttt{step()} does not
        include the contraction margin. This prevents any downstream consumer
        from auditing lawfulness per-step.

  \item \textbf{No \texttt{PIRTMPolicy} Protocol.}
        The \texttt{goal\_budget} constraint is a named contract with no
        runtime enforcement. \texttt{iterate()} accepts bare objects without
        raising \texttt{TypeError}.

  \item \textbf{Unused \texttt{project\_ball()} in \texttt{projection.py}.}
        The function is implemented and documented as appropriate for
        high-dimensional state spaces but is never called by \texttt{iterate()}.

  \item \textbf{Widening spec-to-code gap.}
        60+ spec documents exist in \texttt{docs/specs/} against fewer than
        10 bound code modules. The three largest unbound specs are
        \texttt{ZM\_Meta\_Relativity.md},
        \texttt{Spin\_Foam\_Microfoundations.md}, and
        \texttt{ALE\_GFT.md}.
\end{enumerate}

%% ═══════════════════════════════════════════════════════════════════════════
\section{The ZM Operator Stack: Disclosed Design}
\label{sec:zm-stack}

\subsection{Motivation: Binding Proximity}

The ZM spec is not a separate physics layer --- it is the algebraic completion
of what PIRTM partially implements. The decision criterion for binding ZM first
(ahead of \texttt{Spin\_Foam\_Microfoundations}) is \emph{binding proximity}:
ZM has existing code surface to attach to in \texttt{pirtm/sigma/} and
\texttt{pirtm/core/}, whereas spin foam requires introducing entirely new
module families.

\subsection{ADR-ZM-001: Module Design Decision}

\textbf{Decision:} \texttt{pirtm/zm/operator\_stack.py} \emph{composes from}
existing modules \texttt{pirtm/sigma/} (for kernel weighting) and
\texttt{pirtm/spectral/} (for contractivity margin). It does not own
$D_\sigma$ construction independently, preventing duplication with
\texttt{pirtm/sigma/asymmetric\_kernel.py}.

\textbf{Co-dependency resolution:} \texttt{pirtm/certify.py} and
\texttt{pirtm/zm/operator\_stack.py} must be committed in the same atomic
commit. \texttt{certify.py} does not import from \texttt{pirtm/zm/}
(avoiding circular dependency); \texttt{ZMOperatorStack.certify()} calls
\texttt{certify\_state()} from \texttt{pirtm.certify}.

\subsection{Full Disclosed Module: \texttt{pirtm/zm/operator\_stack.py}}

\begin{lstlisting}[caption={pirtm/zm/operator\_stack.py --- Disclosed stub design},
                   label={lst:zm-stack}]
"""
pirtm/zm/operator_stack.py

Stub binding: ZM_Meta_Relativity.md -> PIRTM operator stack.
Finite-prime truncation of universal operator U = A + C + E.
Full certification via GapLB/SlopeUB stubs raise NotImplementedError
until certify.py accepts (b_p, L_p, w_p) arguments.

ADR: docs/adr/ADR-ZM-001-operator-stack-binding.md
Prior Art Disclosure: Multiplicity Foundation, April 25, 2026.
"""
import numpy as np
from typing import List, Dict, Tuple, Optional


class ZMOperatorStack:
    """
    Finite-prime truncation of
        U = A (x) I (x) I  +  I (x) C (x) I  +  I (x) I (x) E
    on H = ell^2(P) (x) L^2(R) (x) C^d.

    Binding surface for ZM_Meta_Relativity.md s.6.2 Certification Protocol.
    Composes from pirtm.sigma (kernel) and pirtm.spectral (contractivity).
    """

    def __init__(
        self,
        primes: List[int],
        sigma: float = 0.8,
        alpha: float = 0.8,
    ) -> None:
        self.primes = primes
        self.sigma = sigma       # D_sigma diagonal: p^{-sigma}
        self.alpha = alpha       # K HS kernel: p^{-alpha} q^{-alpha} h(.)

        # alpha -> b_p mapping (ZM spec s.2.3):
        #   b_p ~ p^{-alpha} * ||h||_{L^2}
        # Set by build_prime_block(); read by certify().
        self._b_map: Dict[int, float] = {}

        self._A: Optional[np.ndarray] = None
        self._C_range: Optional[Tuple[float, float]] = None

    # ------------------------------------------------------------------
    def build_prime_block(
        self,
        h_fn: Optional[callable] = None
    ) -> np.ndarray:
        """
        Build A = D_sigma + K for finite prime set.
        Delegates D_sigma to pirtm.sigma.asymmetric_kernel.
        Populates self._b_map: b_p = p^{-alpha} * ||h||_{L^2}.

        ADR-ZM-001: implement before certify.py Day 2 fix.
        """
        raise NotImplementedError(
            "ADR-ZM-001: compose from pirtm.sigma before certify.py Day 2"
        )

    def multiplier_range(
        self,
        a0: float,
        a_p: Dict[int, float],
    ) -> Tuple[float, float]:
        """
        Compute [m_min, m_max] for C = F^{-1} M_m F.
        Prerequisite for GapLB (ZM s.6.2 Step 2).
        """
        raise NotImplementedError(
            "ADR-ZM-001: implement before certify.py Day 2 fix"
        )

    def gap_lower_bound(self, b: float, w_budget: float) -> float:
        """
        GapLB >= delta_S - 2 * b * w_budget  (Theorem 6, ZM s.3.2).
        Contractivity margin formula for certify_state() output.

        Wires to certify_state() after Day 2.
        """
        raise NotImplementedError(
            "ADR-ZM-001: wires to certify_state() after Day 2"
        )

    def slope_upper_bound(self, L: float, w_budget: float) -> float:
        """SlopeUB = L * w_budget  (Theorem 6)."""
        raise NotImplementedError("ADR-ZM-001")

    def certify(
        self,
        b_p: Dict[int, float],
        L_p: Dict[int, float],
        w_p: Dict[int, float],
        gamma_min: float,
    ) -> bool:
        """
        Full certification per ZM s.6.2.
        Steps:
          1. Verify HS condition (self._A must be built).
          2. Check multiplier range.
          3. Compute GapLB / SlopeUB.
          4. Return GapLB >= gamma_min.

        Calls pirtm.certify.certify_state() -- co-committed Day 2.
        Raises NotImplementedError until build_prime_block() implemented.
        """
        raise NotImplementedError(
            "ADR-ZM-001: stub -- fails certify test by design"
        )
\end{lstlisting}

\subsection{Full Disclosed Test: \texttt{pirtm/tests/zm/test\_zm\_operator\_stack.py}}

\begin{lstlisting}[caption={Specification test for ZMOperatorStack --- known-failing ADR anchor},
                   label={lst:zm-test}]
"""
pirtm/tests/zm/test_zm_operator_stack.py

ADR-ZM-001 anchor: specification tests, not regression tests.
All tests are expected to fail (NotImplementedError) until
build_prime_block() is implemented after Day 2 certify.py commit.

Prior Art Disclosure: Multiplicity Foundation, April 25, 2026.
"""
import pytest
from pirtm.zm.operator_stack import ZMOperatorStack

FIRST_10_PRIMES = [2, 3, 5, 7, 11, 13, 17, 19, 23, 29]


def test_certify_raises_not_implemented():
    """ADR-ZM-001 anchor: stub raises NotImplementedError."""
    stack = ZMOperatorStack(primes=FIRST_10_PRIMES, sigma=0.8, alpha=0.8)
    with pytest.raises(NotImplementedError):
        stack.certify(b_p={}, L_p={}, w_p={}, gamma_min=0.1)


def test_gap_lower_bound_formula():
    """
    Specification: GapLB(b=1.0, budget=0.1) >= delta_S - 0.2.
    Currently fails -- specification test, not regression test.
    """
    stack = ZMOperatorStack(primes=FIRST_10_PRIMES)
    with pytest.raises(NotImplementedError):
        stack.gap_lower_bound(b=1.0, w_budget=0.1)


def test_b_map_alpha_relationship():
    """
    Specification: after build_prime_block(), b_p ~ p^{-alpha} * ||h||.
    Currently fails -- documents the alpha -> b_p binding contract.
    """
    stack = ZMOperatorStack(primes=FIRST_10_PRIMES, alpha=0.8)
    with pytest.raises(NotImplementedError):
        stack.build_prime_block()
    # Post-implementation assertion (documented here as specification):
    # for p in FIRST_10_PRIMES:
    #     assert stack._b_map[p] == pytest.approx(p**(-0.8), rel=1e-6)


def test_multiplier_range_bounds():
    """
    Specification: multiplier_range returns (m_min, m_max) with
    m_min <= m_max. Currently stub.
    """
    stack = ZMOperatorStack(primes=FIRST_10_PRIMES)
    with pytest.raises(NotImplementedError):
        stack.multiplier_range(a0=1.0, a_p={p: 0.1 for p in FIRST_10_PRIMES})
\end{lstlisting}

%% ═══════════════════════════════════════════════════════════════════════════
\section{The PIRTMPolicy Protocol: Disclosed Design}
\label{sec:policy}

\begin{lstlisting}[caption={\texttt{pirtm/policy.py} --- PIRTMPolicy Protocol stub},
                   label={lst:policy}]
"""
pirtm/policy.py

PIRTMPolicy Protocol: runtime enforcement of the lawfulness inequality
    sum_k eps_k < (1 - gamma_i) * ||P_i||^{-1}
inside iterate(). Bare objects without goal_budget raise TypeError.

Prior Art Disclosure: Multiplicity Foundation, April 25, 2026.
"""
from typing import Protocol, runtime_checkable


@runtime_checkable
class PIRTMPolicy(Protocol):
    """
    Typed protocol for PIRTM iterate() enforcement.

    Any object passed to iterate() as policy= must implement:
      - goal_budget: float   (max perturbation budget)
      - gamma_min: float     (minimum contractivity parameter)
      - operator_norm_bound: float  (upper bound on ||P_i||)

    Violation of the lawfulness inequality raises ValueError inside iterate().
    Absence of this interface raises TypeError at iterate() entry.
    """
    goal_budget: float
    gamma_min: float
    operator_norm_bound: float


def enforce_lawfulness(
    policy: PIRTMPolicy,
    epsilon_sum: float,
) -> None:
    """
    Raise ValueError if lawfulness inequality is violated:
        epsilon_sum >= (1 - gamma_min) * operator_norm_bound^{-1}
    """
    threshold = (1.0 - policy.gamma_min) / policy.operator_norm_bound
    if epsilon_sum >= threshold:
        raise ValueError(
            f"Lawfulness violated: sum(eps_k)={epsilon_sum:.6f} "
            f">= threshold={threshold:.6f}. "
            f"Reduce goal_budget or increase gamma headroom."
        )
\end{lstlisting}

%% ═══════════════════════════════════════════════════════════════════════════
\section{The Five-Gate Ordered Commit Sequence}
\label{sec:gates}

The following ordered sequence constitutes a novel software engineering
artifact for provably-correct PIRTM deployment. The ordering is
dependency-driven: each gate unblocks the next.

\begin{longtable}{cllp{4.5cm}}
\toprule
\textbf{Gate} & \textbf{File} & \textbf{Metric} & \textbf{Dependency} \\
\midrule
1 & \texttt{pirtm/core/recurrence.py} &
    \texttt{StepInfo.margin} present on 100\% of step calls &
    None (standalone) \\[4pt]

2 & \texttt{pirtm/certify.py} \newline
    \texttt{pirtm/zm/\_\_init\_\_.py} \newline
    \texttt{pirtm/zm/operator\_stack.py} &
    \texttt{from pirtm.certify import certify\_state} succeeds;
    single atomic commit &
    Gate 1 (StepInfo.margin available) \\[4pt]

3 & \texttt{pirtm/tests/test\_supermodule\_spectral.py} \newline
    \texttt{pirtm/tests/zm/test\_zm\_operator\_stack.py} &
    Both committed as known-failing spec tests;
    \texttt{pytest} exits with known-failure &
    Gate 2 (import targets exist) \\[4pt]

4 & \texttt{pirtm/policy.py} &
    \texttt{iterate()} raises \texttt{TypeError} on bare object;
    raises \texttt{ValueError} on lawfulness violation &
    Gate 2 (certify\_state computable) \\[4pt]

5 & \texttt{docs/adr/ADR-ZM-001.md} \newline
    \texttt{docs/adr/ADR-supermodule.md} &
    ADRs reference test anchors from Gate 3;
    Section 3 of ADR-ZM-001 states compose-vs-own decision &
    Gate 3 (test anchors exist) \\
\bottomrule
\caption{Five-gate ordered commit sequence for PIRTM provable-correctness deployment.}
\end{longtable}

\textbf{Co-dependency note (Gate 2):} \texttt{pirtm/certify.py} and
\texttt{pirtm/zm/operator\_stack.py} are co-dependent --- \texttt{certify.py}
needs the parameter names that ZM defines, and \texttt{ZMOperatorStack.certify()}
calls \texttt{certify\_state()}. Splitting Gate 2 into two commits creates an
import-broken window. The single-commit resolution is the disclosed design.

%% ═══════════════════════════════════════════════════════════════════════════
\section{Spin Foam as Prime-Gated Hilbert Space}
\label{sec:spin-foam}

The identification of spin foam microfoundations as an instantiation of the
ZM prime-gated Hilbert space $\ell^2(\mathbb{P})$ is a novel contribution
of this framework. In the spin foam context, area eigenvalues of the area
operator $\hat{A}$ take the form:
\[
  a_j \;=\; \ell_P^2 \sqrt{j(j+1)},
  \quad j \in \tfrac{1}{2}\mathbb{Z}_{\geq 0},
\]
where $\ell_P$ is the Planck length. The \emph{disclosed integration claim}
is: the spin labels $j$ play the role of prime weights in $\ell^2(\mathbb{P})$
under the identification
\[
  p \;\longleftrightarrow\; j(j+1),
  \quad
  p^{-\sigma} \;\longleftrightarrow\; \ell_P^{-2}\, a_j,
\]
enabling spin network states to be processed through the ZM operator stack as
a specific instantiation of prime-gated Hilbert space structure. Spin foam
binding (\texttt{pirtm/spin\_foam/}) is the designated Week 3 target,
following ZM stub establishment at Gate 2.

%% ═══════════════════════════════════════════════════════════════════════════
\section{AL-GFT: Gravity as Complex Open System}
\label{sec:algft}

The Asymmetric Langevin Gravity Field Theory (AL-GFT) specification
(\texttt{AL-GFT\_CEQG-RG-Langevin.md}, 124,835 bytes; \texttt{ALE\_GFT.md},
170,419 bytes) extends the Meta-Relativity framework to a stochastic
treatment of gravitational degrees of freedom. The key disclosed elements are:

\begin{itemize}[leftmargin=*]
  \item A \emph{Coarse-grained Effective Quantum Gravity} (CEQG) layer in which
        the renormalization group (RG) flow is governed by a Langevin equation
        with prime-indexed noise kernel;
  \item The interpretation of the RG $\beta$-function as the contractivity
        slope $\operatorname{SlopeUB}$ from Theorem~\ref{thm:gaplb};
  \item A \emph{regime dossier system} (Boundary-Spectral Governance,
        \texttt{Boundary-spectral Governance --- Regime Dossier System v1.md})
        classifying gravitational regimes by their spectral gap position
        relative to $\operatorname{GapLB}$.
\end{itemize}

%% ═══════════════════════════════════════════════════════════════════════════
\section{Prior Art Disclosure Summary}
\label{sec:prior-art}

The following table summarizes all novel elements disclosed by this publication
as of April 25, 2026:

\begin{longtable}{p{4.2cm}p{3.5cm}p{5.5cm}}
\toprule
\textbf{Disclosed Element} & \textbf{Location} & \textbf{Claims Prevented} \\
\midrule
ZM universal operator $U = A \otimes I + I \otimes C \otimes I + I \otimes I \otimes E$
& \S\ref{sec:zm}, Def.~\ref{def:zm-universal}
& Patent on tensor-product prime-indexed operator for physics simulation \\[4pt]

GapLB/SlopeUB certification bounds (Theorem~\ref{thm:gaplb})
& \S\ref{sec:cert}
& Patent on contractivity certification via per-prime Lipschitz bounds \\[4pt]

$\alpha \mapsto b_p$ binding ($b_p \sim p^{-\alpha}\|h\|_{L^2}$)
& \S\ref{sec:cert}, Remark
& Patent on HS kernel parameter to per-prime bound mapping \\[4pt]

PIRTMPolicy Protocol + \texttt{enforce\_lawfulness()}
& \S\ref{sec:policy}, Listing~\ref{lst:policy}
& Patent on typed protocol for operator-norm-bounded policy enforcement \\[4pt]

ZMOperatorStack stub architecture + ADR co-dependency pattern
& \S\ref{sec:zm-stack}, Listings \ref{lst:zm-stack}--\ref{lst:zm-test}
& Patent on stub-first ADR-governed module binding for operator stacks \\[4pt]

Five-gate ordered commit sequence
& \S\ref{sec:gates}
& Patent on dependency-ordered commit sequencing for provable-correctness deployment \\[4pt]

Spin foam as $\ell^2(\mathbb{P})$ instantiation
& \S\ref{sec:spin-foam}
& Patent on prime-weight identification of spin foam area eigenvalues \\[4pt]

AL-GFT RG $\beta$-function as SlopeUB
& \S\ref{sec:algft}
& Patent on renormalization group / contractivity-slope identification \\[4pt]

\texttt{StepInfo.margin} field design
& \S\ref{sec:arch}
& Patent on per-step contractivity margin in operator iteration dataclass \\[4pt]

\texttt{certify\_state(b\_p, L\_p, w\_p)} signature
& \S\ref{sec:cert}, \S\ref{sec:arch}
& Patent on per-prime argument certification function signature \\
\bottomrule
\caption{Prior art disclosure summary.}
\end{longtable}

%% ═══════════════════════════════════════════════════════════════════════════
\section{Recommended Next Steps}
\label{sec:next}

\begin{enumerate}[leftmargin=*]
  \item \textbf{Submit this document} to a timestamped public repository
        (e.g., arXiv cs.MS or math.OA, OSF, or IP.com Technical Disclosure
        Commons) immediately to establish a public disclosure date with
        independent timestamp.

  \item \textbf{Engage an IP attorney} specializing in mathematical software
        and computer-implemented inventions to evaluate whether any elements
        of the ZM operator construction or GapLB certification protocol
        rise to the level of patentable subject matter under
        35 U.S.C.~\S\,101 post-\textit{Alice}, and to assess freedom-to-operate
        against existing filings in operator theory software.

  \item \textbf{File a provisional patent application} (if advised by counsel)
        within 12 months of this disclosure date, preserving patent rights
        for elements not yet in the public domain.

  \item \textbf{Complete Gate 2} (atomic commit of \texttt{pirtm/certify.py}
        and \texttt{pirtm/zm/}) to create immutable git evidence of
        implementation date.

  \item \textbf{Register the copyright} for the spec corpus
        (\texttt{docs/specs/}, 60+ documents) with the U.S. Copyright Office
        as a compilation, establishing a public registration date.
\end{enumerate}


%% ─── APPENDIX ─────────────────────────────────────────────────────────────
\appendix

\section{Finite Prime Truncation: SageMath Reference Implementation}
\label{app:sage}

The following is a directly translatable Python reference for the finite-prime
truncation of $A = D_\sigma + K$ (from ZM spec \S6.4), included as part of
the prior art disclosure:

\begin{lstlisting}[caption={Finite prime truncation reference (translatable from ZM s.6.4)},
                   label={lst:sage}]
import numpy as np

def h_kernel(x: float, scale: float = 1.0) -> float:
    """Gaussian kernel h(log p - log q)."""
    return np.exp(-0.5 * (x / scale) ** 2)

def build_prime_block_reference(
    primes: list,
    sigma: float = 0.8,
    alpha: float = 0.8,
    h_scale: float = 1.0,
) -> tuple:
    """
    Reference implementation: A = D_sigma + K for finite prime set.
    Returns (A, b_map) where b_map[p] = p^{-alpha} * ||h||_approx.
    NOT the production implementation -- see pirtm/zm/operator_stack.py.
    Prior Art Disclosure: Multiplicity Foundation, April 25, 2026.
    """
    n = len(primes)
    # D_sigma: diagonal p^{-sigma}
    D = np.diag([p ** (-sigma) for p in primes])
    # K: Hilbert-Schmidt off-diagonal
    K = np.zeros((n, n))
    for i, p in enumerate(primes):
        for j, q in enumerate(primes):
            K[i, j] = (p ** (-alpha)) * (q ** (-alpha)) * \
                       h_kernel(np.log(p) - np.log(q), scale=h_scale)
    A = D + K
    # b_map: alpha -> b_p binding (ZM s.2.3)
    h_norm_approx = np.sqrt(np.sum(
        h_kernel(np.log(p) - np.log(q), scale=h_scale) ** 2
        for p in primes for q in primes
    ) / n)
    b_map = {p: p ** (-alpha) * h_norm_approx for p in primes}
    return A, b_map

def gap_lower_bound_reference(
    delta_S: float,
    b_map: dict,
    w_p: dict,
) -> float:
    """
    GapLB >= delta_S - 2 * sum_p |w_p| * b_p  (Theorem 6).
    Reference implementation for prior art disclosure.
    """
    correction = 2.0 * sum(abs(w_p.get(p, 0.0)) * b_map[p] for p in b_map)
    return delta_S - correction

# Example evaluation with first 10 primes
if __name__ == "__main__":
    primes = [2, 3, 5, 7, 11, 13, 17, 19, 23, 29]
    A, b_map = build_prime_block_reference(primes, sigma=0.8, alpha=0.8)
    w_p = {p: 0.05 for p in primes}
    delta_S = 0.45  # example spectral gap
    gaplb = gap_lower_bound_reference(delta_S, b_map, w_p)
    print(f"||A||_op approx: {np.linalg.norm(A, ord=2):.6f}")
    print(f"GapLB: {gaplb:.6f}")
    print(f"Lawful: {gaplb > 0.1}")
\end{lstlisting}

\section{Spec Corpus Inventory (Selected)}
\label{app:specs}

\begin{longtable}{lrl}
\toprule
\textbf{Specification File} & \textbf{Size (bytes)} & \textbf{Binding Status} \\
\midrule
\texttt{ALE\_GFT.md} & 170,419 & Unbound \\
\texttt{AL-GFT\_CEQG-RG-Langevin.md} & 124,835 & Unbound \\
\texttt{ZM\_Meta\_Relativity.md} & $\sim$77,000 & Gate 2 target \\
\texttt{Spin\_Foam\_Microfoundations.md} & $\sim$75,000 & Week 3 target \\
\texttt{Meta-Ensembles.md} & 30,880 & Unbound \\
\texttt{Aspectual\_Counting\_Framework.md} & 30,486 & Unbound \\
\texttt{Meta-relativity.md} & 25,524 & Partial \\
\texttt{Meta-Relativity Framework: Operational Manual} & 17,010 & Partial \\
\texttt{Atomic\_Multiplicity.md} & 14,780 & Unbound \\
\texttt{Multiplicity-inflected Prime Theory.md} & 10,248 & Partial \\
\bottomrule
\caption{Selected spec corpus entries and binding status (April 2026).}
\end{longtable}

============================================================
\section{Notation and Standing Assumptions}
\label{app:notation}
% ============================================================

Throughout this appendix $\PP = \{2,3,5,7,11,\ldots\}$ denotes the set of
rational primes. We write $\lp$ for the separable Hilbert space of complex
square-summable sequences indexed by $\PP$, with inner product
\[
  \ip{f}{g}_{\lp} \;=\; \sum_{p \in \PP} f(p)\,\overline{g(p)}.
\]
The full three-component state space is
\begin{equation}
  \HH \;=\; \lp \;\otimes\; \Lt \;\otimes\; \Cd,
  \label{eq:hilbert}
\end{equation}
where $L^2(\RR)$ carries the standard Lebesgue measure and $\Cd$ is a
finite-dimensional environment of dimension $d \geq 1$.

\medskip\noindent\textbf{Standing Assumptions.}
\begin{enumerate}[label=\textbf{(SA\arabic*)}]
  \item\label{sa:sigma} $\sigma > \tfrac{1}{2}$. (Ensures absolute convergence
        of $\sum_p p^{-2\sigma}$ by the prime zeta function bound.)
  \item\label{sa:alpha} $\alpha > \tfrac{1}{2}$. (Ensures the kernel $K$ is
        Hilbert-Schmidt; see \cref{app:hs}.)
  \item\label{sa:h} $h \in L^2(\RR)$ with $\norm{h}_{L^2} < \infty$.
  \item\label{sa:m} The Fourier multiplier symbol $m \in L^\infty(\RR)$ is
        real-valued with $\essran(m) \subseteq [m_{\min}, m_{\max}]$,
        $-\infty < m_{\min} \leq m_{\max} < \infty$.
  \item\label{sa:E} $E \in M_d(\CC)$ is normal with $\norm{E}_{\op} < \infty$.
\end{enumerate}

We write $\norm{\,\cdot\,}_{\op}$ for the operator norm, $\norm{\,\cdot\,}_{\HS}$
for the Hilbert-Schmidt norm, and $\norm{\,\cdot\,}_{p}$ for the Schatten
$p$-norm. The symbol $\lesssim$ denotes inequality up to a universal constant.

% ============================================================
\section{Hilbert-Schmidt Norm of the Off-Diagonal Kernel \texorpdfstring{$K$}{K}}
\label{app:hs}
% ============================================================

\begin{appdefn}[Off-diagonal kernel]
\label{defn:K}
For $p,q \in \PP$ define
\[
  K(p,q) \;=\; p^{-\alpha}\, q^{-\alpha}\, h\!\left(\log p - \log q\right),
\]
and regard $K$ as the integral kernel of the operator
$K \colon \lp \to \lp$, $(Kf)(p) = \sum_{q \in \PP} K(p,q)\,f(q)$.
\end{appdefn}

\begin{apptheorem}[Hilbert-Schmidt bound for $K$]
\label{thm:hs}
Under Assumptions \ref{sa:alpha} and \ref{sa:h},
\[
  \norm{K}_{\HS}^2
  \;=\; \sum_{p,q \in \PP} \abs{K(p,q)}^2
  \;\leq\; \zeta_\PP(2\alpha)^2 \cdot \norm{h}_{L^2}^2
  \;<\; \infty,
\]
where $\zeta_\PP(s) = \sum_{p \in \PP} p^{-s}$ is the prime zeta function,
absolutely convergent for $\Re(s) > 1$.
\end{apptheorem}

\begin{proof}
Compute the Hilbert-Schmidt norm directly:
\begin{align}
  \norm{K}_{\HS}^2
  &= \sum_{p,q \in \PP} \abs{K(p,q)}^2
   = \sum_{p,q \in \PP} p^{-2\alpha}\, q^{-2\alpha}\,
     \abs{h(\log p - \log q)}^2.
  \label{eq:hs-sum}
\end{align}
Write $x_{p,q} = \log p - \log q$. For fixed $p$, the map $q \mapsto x_{p,q}$
is injective on $\PP$, so we may bound by extending the sum over all
$(p,q) \in \PP \times \PP$. Factor:
\begin{align}
  \sum_{p,q \in \PP}
    p^{-2\alpha} q^{-2\alpha} \abs{h(\log p - \log q)}^2
  &\leq
  \left(\sum_{p \in \PP} p^{-2\alpha}\right)
  \cdot
  \sup_{p}\!\left[\sum_{q \in \PP} q^{-2\alpha}\,
    \abs{h(\log p - \log q)}^2\right].
  \label{eq:hs-factor}
\end{align}
For the inner sum, change variable $u = \log q$ and bound the discrete sum
by the continuous integral (using the mean-value theorem for the prime-counting
function $\pi(x) \sim x/\log x$):
\begin{align}
  \sum_{q \in \PP} q^{-2\alpha} \abs{h(\log p - \log q)}^2
  &\lesssim \int_0^\infty e^{-2\alpha u} \abs{h(\log p - u)}^2\,
    \frac{e^u\,du}{u} \notag\\
  &\leq \norm{h}_{L^2}^2 \cdot
    \sup_{u > 0}\!\left[\frac{e^{(1-2\alpha)u}}{u}\right].
  \label{eq:inner-integral}
\end{align}
For $\alpha > \tfrac{1}{2}$ the exponent $1-2\alpha < 0$, so
$\sup_{u>0} e^{(1-2\alpha)u}/u < \infty$.
Returning to \eqref{eq:hs-factor} and using
$\sum_{p \in \PP} p^{-2\alpha} = \zeta_\PP(2\alpha) < \infty$ for
$2\alpha > 1$, we obtain:
\[
  \norm{K}_{\HS}^2
  \;\leq\; \zeta_\PP(2\alpha)^2\, \norm{h}_{L^2}^2
  \;<\; \infty. \qedhere
\]
\end{proof}

\begin{appcorollary}[Compactness of $K$]
\label{cor:compact}
Under Assumptions \ref{sa:alpha}--\ref{sa:h}, $K$ is a compact operator on
$\lp$, and in particular $\spec_{\mathrm{ess}}(A) = \spec_{\mathrm{ess}}(D_\sigma)$.
\end{appcorollary}

\begin{proof}
Every Hilbert-Schmidt operator is compact \cite[Thm.~VI.22]{reed_simon}.
By the Weyl theorem, compact perturbations preserve the essential spectrum,
giving $\spec_{\mathrm{ess}}(D_\sigma + K) = \spec_{\mathrm{ess}}(D_\sigma)$.
\end{proof}

% ============================================================
\section{Operator Norm Bound for the Prime Block \texorpdfstring{$A$}{A}}
\label{app:prime-block}
% ============================================================

\begin{apptheorem}[Operator norm of $A$]
\label{thm:A-norm}
Under Assumptions \ref{sa:sigma}, \ref{sa:alpha}, \ref{sa:h},
\[
  \norm{A}_{\op}
  \;\leq\;
  \norm{D_\sigma}_{\op} + \norm{K}_{\op}
  \;\leq\;
  2^{-\sigma} + \zeta_\PP(2\alpha)\,\norm{h}_{L^2},
\]
where $2^{-\sigma}$ is the spectral radius of $D_\sigma$ (attained at $p=2$).
\end{apptheorem}

\begin{proof}
By the triangle inequality for the operator norm:
\[
  \norm{A}_{\op} \leq \norm{D_\sigma}_{\op} + \norm{K}_{\op}.
\]
\textit{Bound for $D_\sigma$.}
$D_\sigma$ is diagonal with entries $p^{-\sigma}$, so
$\norm{D_\sigma}_{\op} = \sup_{p \in \PP} p^{-\sigma} = 2^{-\sigma}$.

\textit{Bound for $K$.}
For any Hilbert-Schmidt operator, $\norm{K}_{\op} \leq \norm{K}_{\HS}$.
Applying \cref{thm:hs}:
\[
  \norm{K}_{\op} \leq \norm{K}_{\HS} \leq \zeta_\PP(2\alpha)\,\norm{h}_{L^2}.
\]
Combining yields the stated bound.
\end{proof}

\begin{appremark}
For concrete parameters $\sigma = 0.8$, $\alpha = 0.8$, $\norm{h}_{L^2} = 1$:
\[
  \norm{A}_{\op} \;\leq\; 2^{-0.8} + \zeta_\PP(1.6)\,
  \approx\; 0.574 + 0.452 \;=\; 1.026.
\]
The prime zeta value $\zeta_\PP(1.6) \approx 0.452$ is computed from the
first convergent partial sum over primes up to $10^6$.
\end{appremark}

% ============================================================
\section{Spectral Gap of \texorpdfstring{$A$}{A}}
\label{app:specgap}
% ============================================================

\begin{appdefn}[Phase-parameterized spectral gap]
\label{defn:specgap}
For $\theta \in [0, 2\pi)$ define the \emph{phase-$\theta$ spectral gap} of
$A$ as
\[
  \sgap(\theta)
  \;=\;
  \inf\!\left\{
    \abs{e^{i\theta}\lambda - \mu}
    \;\colon\;
    \lambda \in \spec(A),\; \mu \in \spec(A),\; \lambda \neq \mu
  \right\}.
\]
The \emph{global spectral gap} is $\sgap^* = \inf_\theta \sgap(\theta)$.
\end{appdefn}

\begin{applemma}[Lower bound on $\sgap^*$ from $D_\sigma$]
\label{lem:sgap-lower}
If $K = 0$ (pure diagonal), then
\[
  \sgap^*
  \;\geq\;
  \min_{p \neq q \in \PP} \abs{p^{-\sigma} - q^{-\sigma}}
  \;=\;
  3^{-\sigma} - 2^{-\sigma} \cdot \left[1 - (3/2)^{-\sigma}\right]^{-1}
  \;\cdot\; \text{(gap at } p=2, q=3\text{)}.
\]
For $K \neq 0$, by the Bauer-Fike theorem:
\[
  \sgap^*(A) \;\geq\; \sgap^*(D_\sigma) - 2\norm{K}_{\op}.
\]
\end{applemma}

\begin{proof}
For the diagonal case, eigenvalues are $\{p^{-\sigma}\}_{p \in \PP}$, which
are a strictly decreasing sequence. The minimum gap is between consecutive
primes $2$ and $3$: $3^{-\sigma} - 2^{-\sigma}$... wait, both are positive
and $2^{-\sigma} > 3^{-\sigma}$, so
$\sgap^*(D_\sigma) = 2^{-\sigma} - 3^{-\sigma} > 0$ for all finite $\sigma$.

For the perturbed case, the Bauer-Fike theorem states that for any eigenvalue
$\mu$ of $A = D_\sigma + K$ there exists an eigenvalue $\lambda$ of
$D_\sigma$ such that $\abs{\mu - \lambda} \leq \norm{K}_{\op}$. Hence
each eigenvalue of $A$ lies within an $\norm{K}_{\op}$-ball around some
eigenvalue of $D_\sigma$. The minimum separation between distinct eigenvalue
clusters is at least $\sgap^*(D_\sigma) - 2\norm{K}_{\op}$, provided
$\norm{K}_{\op} < \tfrac{1}{2}\sgap^*(D_\sigma)$.
\end{proof}

\begin{appcorollary}[Non-degeneracy condition]
\label{cor:nondegen}
$A$ has a positive spectral gap, i.e., $\sgap^*(A) > 0$, whenever
\[
  \norm{K}_{\op}
  \;<\;
  \frac{1}{2}(2^{-\sigma} - 3^{-\sigma}).
\]
For $\sigma = 0.8$: $\tfrac{1}{2}(2^{-0.8} - 3^{-0.8}) \approx 0.132$.
\end{appcorollary}

% ============================================================
\section{Norm of the Universal Operator \texorpdfstring{$U$}{U}}
\label{app:U-norm}
% ============================================================

\begin{apptheorem}[Operator norm of $U$]
\label{thm:U-norm}
Under Assumptions \ref{sa:sigma}--\ref{sa:E}, the universal operator
\[
  U = A \otimes I \otimes I + I \otimes C \otimes I + I \otimes I \otimes E
\]
is a bounded self-adjoint operator on $\HH$ with
\[
  \norm{U}_{\op}
  \;\leq\;
  \norm{A}_{\op} + \norm{C}_{\op} + \norm{E}_{\op}
  \;\leq\;
  \left(2^{-\sigma} + \zeta_\PP(2\alpha)\norm{h}_{L^2}\right)
  + \max(\abs{m_{\min}},\abs{m_{\max}})
  + \norm{E}_{\op}.
\]
\end{apptheorem}

\begin{proof}
By the triangle inequality for operator norms on tensor products:
\[
  \norm{A \otimes I \otimes I}_{\op} = \norm{A}_{\op}\cdot\norm{I}_{\op}^2
  = \norm{A}_{\op},
\]
and analogously for the other two summands. Hence:
\[
  \norm{U}_{\op}
  \leq \norm{A}_{\op} + \norm{C}_{\op} + \norm{E}_{\op}.
\]
\textit{Bound for $C$.}
$C = \mathcal{F}^{-1} M_m \mathcal{F}$; since $\mathcal{F}$ is unitary,
$\norm{C}_{\op} = \norm{M_m}_{\op} = \norm{m}_{L^\infty} = \max(\abs{m_{\min}},\abs{m_{\max}})$.

\textit{Bound for $A$.} Applied from \cref{thm:A-norm}.

Combining gives the stated bound.
\end{proof}

\begin{appcorollary}[Sufficient condition for $\norm{U}_{\op} < 1$]
\label{cor:contraction}
$U$ is a contraction ($\norm{U}_{\op} \leq 1$) if
\[
  2^{-\sigma} + \zeta_\PP(2\alpha)\norm{h}_{L^2}
  + \max(\abs{m_{\min}},\abs{m_{\max}})
  + \norm{E}_{\op} \;\leq\; 1.
\]
\end{appcorollary}

% ============================================================
\section{The Recursive Operator \texorpdfstring{$\Xi(t)$}{Xi(t)}: Convergence and Contractivity}
\label{app:xi}
% ============================================================

\begin{appdefn}[Recursive operator semigroup]
\label{defn:xi}
Let $T \in \mathcal{B}(\HH)$ with $\norm{T}_{\op} \leq 1$. Define the
\emph{recursive operator} by the power series
\[
  \Xi(t) \;=\; \sum_{k=0}^{\infty} \alpha_k(t)\, T^k,
\]
where $\{\alpha_k(t)\}_{k \geq 0}$ are real-valued coefficients satisfying
for all $t \geq 0$:
\begin{enumerate}[label=(\roman*)]
  \item $\alpha_k(t) \geq 0$ for all $k$,
  \item $\sum_{k=0}^\infty \alpha_k(t) = 1$ (convex combination),
  \item $\alpha_0(0) = 1$ (initial condition: $\Xi(0) = I$).
\end{enumerate}
\end{appdefn}

\begin{apptheorem}[Well-posedness and contractivity of $\Xi(t)$]
\label{thm:xi-contraction}
Under the conditions of \cref{defn:xi} and $\norm{T}_{\op} \leq \gamma < 1$:
\begin{enumerate}[label=(\roman*)]
  \item The series $\Xi(t) = \sum_k \alpha_k(t) T^k$ converges in operator
        norm for all $t \geq 0$.
  \item $\norm{\Xi(t)}_{\op} \leq 1$ for all $t \geq 0$.
  \item $\Xi(t)$ is a $\gamma$-contraction: for all $\psi, \phi \in \HH$,
        \[
          \norm{\Xi(t)\psi - \Xi(t)\phi}
          \;\leq\; \gamma_{\mathrm{eff}}(t)\,\norm{\psi - \phi},
        \]
        where
        \[
          \gamma_{\mathrm{eff}}(t)
          \;=\;
          \sum_{k=1}^{\infty} \alpha_k(t)\, \gamma^k
          \;\leq\;
          \gamma \sum_{k=1}^{\infty} \alpha_k(t)
          \;=\;
          \gamma\,(1 - \alpha_0(t)).
        \]
\end{enumerate}
\end{apptheorem}

\begin{proof}
\textit{(i) Convergence.}
Since $\norm{T^k}_{\op} \leq \norm{T}_{\op}^k \leq 1$:
\[
  \norm{\Xi(t)}_{\op}
  \leq \sum_{k=0}^\infty \alpha_k(t)\,\norm{T^k}_{\op}
  \leq \sum_{k=0}^\infty \alpha_k(t) = 1 < \infty.
\]
The partial sums form a Cauchy sequence in $\mathcal{B}(\HH)$ (Banach), hence
convergent.

\textit{(ii) Norm bound.}
Follows immediately from the estimate above.

\textit{(iii) Contractivity.}
For any $\psi, \phi \in \HH$:
\begin{align}
  \norm{\Xi(t)\psi - \Xi(t)\phi}
  &= \norm*{\sum_{k=0}^\infty \alpha_k(t) T^k (\psi - \phi)} \notag\\
  &\leq \sum_{k=0}^\infty \alpha_k(t)\,\norm{T^k}_{\op}\,\norm{\psi - \phi}
  \notag\\
  &\leq \left[\alpha_0(t) + \sum_{k=1}^\infty \alpha_k(t)\,\gamma^k \right]
    \norm{\psi - \phi}.
\end{align}
Since $\alpha_0(t) \cdot \norm{T^0}_{\op} = \alpha_0(t)$ contributes the
identity component (no contraction), the effective contraction parameter is
\[
  \gamma_{\mathrm{eff}}(t)
  = \alpha_0(t) + \sum_{k=1}^\infty \alpha_k(t)\gamma^k
  \leq \alpha_0(t) + \gamma(1 - \alpha_0(t)).
\]
For $\alpha_0(t) \to 0$ as $t \to \infty$ (mixing condition), we obtain
$\gamma_{\mathrm{eff}}(t) \to \gamma < 1$.
\end{proof}

\begin{appremark}[Relation to \texttt{iterate()}]
The iterate index $k$ in $\sum_k \alpha_k(t) T^k$ corresponds directly to
the loop counter in \texttt{pirtm/core/recurrence.py::iterate()}.
The contraction parameter $\gamma_{\mathrm{eff}}(t)$ is the quantity
that \texttt{StepInfo.margin} must track at each step. Specifically,
\[
  \texttt{StepInfo.margin}
  \;\equiv\;
  1 - \gamma_{\mathrm{eff}}(t_{\mathrm{current}})
  \;=\;
  (1 - \alpha_0(t))(1 - \gamma).
\]
\end{appremark}

% ============================================================
\section{GapLB and SlopeUB: Proof of Theorem 6}
\label{app:gaplb}
% ============================================================

\begin{apptheorem}[GapLB/SlopeUB bounds --- full proof]
\label{thm:gaplb-full}
Let $\{b_p\}_{p \in \PP}$ be per-prime Lipschitz bounds with
$b_p = p^{-\alpha}\norm{h}_{L^2}$. Let $\{L_p\}_{p \in \PP}$ be per-prime
slope bounds with $L = \sup_p L_p < \infty$. Let
$w_{\mathrm{budget}} = \sum_p \abs{w_p}$. Define
\[
  \delta_S(\theta)
  \;=\;
  \inf\!\left\{
    \Re\!\left(e^{-i\theta}\ip{Af}{f}\right)
    \;\colon\; f \in \lp,\; \norm{f} = 1
  \right\},
\]
the phase-$\theta$ numerical range lower bound of $A$. Then:
\begin{align}
  \GapLB(w)
  &\;\geq\;
  \inf_{\theta \in [0,2\pi)}
  \left[\delta_S(\theta) - 2\sum_{p \in \PP} \abs{w_p}\, b_p \right],
  \label{eq:gaplb} \\[6pt]
  \SlopeUB(w)
  &\;=\; L \cdot w_{\mathrm{budget}}.
  \label{eq:slopeub}
\end{align}
\end{apptheorem}

\begin{proof}
\textit{Step 1: Reduction to the prime block.}
The contractivity margin of $U = A \otimes I + I \otimes C \otimes I
+ I \otimes I \otimes E$ on product states
$f \otimes g \otimes v$ decomposes as
\[
  \Re\,\ip{Uf}{f}_\HH = \Re\,\ip{Af}{f}_{\lp}
  + \Re\,\ip{Cg}{g}_{\Lt}
  + \Re\,\ip{Ev}{v}_{\Cd}.
\]
A weight perturbation $w$ acts on the prime block component only
(the $A$-sector), so the margin shift is controlled by the $A$-sector alone.
It suffices to prove the bound for $A$.

\textit{Step 2: Numerical range perturbation.}
The operator $A(w) = A + W$ where $W$ is the weight perturbation operator
with $\norm{W}_{\op} \leq \sum_p \abs{w_p}\, b_p$ (by definition of $b_p$
as the per-prime Lipschitz bound). By the numerical range perturbation lemma
(a standard consequence of the Cauchy-Schwarz inequality):
\[
  \delta_S(\theta, A(w))
  \;\geq\;
  \delta_S(\theta, A) - \norm{W}_{\op}
  \;\geq\;
  \delta_S(\theta) - \sum_p \abs{w_p}\, b_p.
\]

\textit{Step 3: GapLB lower bound.}
The contractivity gap (the infimum of the real part of the spectrum of $-A(w)$
over all unit vectors and all phases $\theta$) is bounded below by:
\[
  \GapLB(w)
  \;\geq\;
  \inf_\theta \left[\delta_S(\theta) - 2\sum_p \abs{w_p}\, b_p\right],
\]
where the factor of $2$ arises from the two-sided spectral perturbation
(both the positive and negative real axis contributions are affected).

\textit{Step 4: SlopeUB.}
The slope of the contractivity bound as a function of $w_{\mathrm{budget}}$
is bounded above by the maximum per-prime Lipschitz constant $L = \sup_p L_p$:
\[
  \abs{\partial_{w_{\mathrm{budget}}} \GapLB(w)}
  \;\leq\;
  2 \sup_p b_p
  \;\leq\;
  2L,
\]
giving $\SlopeUB(w) = L \cdot w_{\mathrm{budget}}$ as the integrated slope
upper bound over the budget interval $[0, w_{\mathrm{budget}}]$.
\end{proof}

\begin{appcorollary}[Certifiability criterion]
\label{cor:cert}
A configuration is certifiably lawful with contractivity parameter
$\gamma_{\min}$ if and only if
\[
  \inf_\theta\!\left[\delta_S(\theta)\right]
  \;>\;
  \gamma_{\min} + 2\sum_p \abs{w_p}\, b_p.
\]
This is the condition evaluated by \texttt{ZMOperatorStack.certify()} and
\texttt{pirtm.certify.certify\_state()}.
\end{appcorollary}

% ============================================================
\section{The PIRTMPolicy Lawfulness Inequality}
\label{app:policy}
% ============================================================

\begin{apptheorem}[PIRTMPolicy lawfulness bound]
\label{thm:policy}
Let $\{P_i\}_{i=1}^n$ be a family of projections on $\HH$ with
$\norm{P_i}_{\op} < \infty$, let $\gamma_i \in (0,1)$ be the contractivity
parameter for module $i$, and let $\{\varepsilon_k\}_{k \geq 1}$ be
step-level perturbation magnitudes. Then the iteration \texttt{iterate()}
preserves lawfulness --- i.e., $\Xi(t)$ remains a contraction after $n$
steps --- if and only if
\begin{equation}
  \sum_{k=1}^n \varepsilon_k
  \;<\;
  (1 - \gamma_i) \cdot \norm{P_i}_{\op}^{-1}
  \label{eq:policy}
\end{equation}
for each active module $i$.
\end{apptheorem}

\begin{proof}
After $n$ steps the accumulated perturbation to the operator norm is bounded
by $\sum_k \varepsilon_k$. For $\Xi(t)$ to remain a $\gamma_i$-contraction
on the range of $P_i$, we require that the perturbed operator
$\tilde T = T + \Delta$ with $\norm{\Delta}_{\op} \leq \sum_k \varepsilon_k$
satisfies
\[
  \norm{\tilde T P_i}_{\op}
  \leq \norm{T P_i}_{\op} + \norm{\Delta P_i}_{\op}
  \leq \gamma_i + \left(\sum_k \varepsilon_k\right)\norm{P_i}_{\op}
  < 1.
\]
Rearranging:
\[
  \left(\sum_k \varepsilon_k\right)\norm{P_i}_{\op}
  < 1 - \gamma_i,
\]
which is exactly \eqref{eq:policy}.
\end{proof}

\begin{appremark}[Runtime enforcement]
\cref{thm:policy} is the mathematical justification for the
\texttt{enforce\_lawfulness()} function in \texttt{pirtm/policy.py}:
\[
  \texttt{threshold}
  = \frac{1 - \texttt{policy.gamma\_min}}{\texttt{policy.operator\_norm\_bound}}
  \equiv
  (1 - \gamma_i) \cdot \norm{P_i}_{\op}^{-1}.
\]
The \texttt{ValueError} raised when
$\texttt{epsilon\_sum} \geq \texttt{threshold}$ is the runtime
manifestation of the violation of \eqref{eq:policy}.
\end{appremark}

% ============================================================
\section{Alpha-to-\texorpdfstring{$b_p$}{b\_p} Binding Map}
\label{app:alpha-bp}
% ============================================================

\begin{appproposition}[$\alpha \mapsto b_p$ explicit form]
\label{prop:alpha-bp}
Under \cref{defn:K}, the per-prime Lipschitz bound $b_p$ induced by the
kernel parameter $\alpha$ is:
\[
  b_p \;=\; p^{-\alpha}\, \norm{h}_{L^2},
\]
and the aggregate budget satisfies
\[
  \sum_p \abs{w_p}\, b_p
  \;=\;
  \norm{h}_{L^2} \sum_p \abs{w_p}\, p^{-\alpha}.
\]
\end{appproposition}

\begin{proof}
The Lipschitz constant of the map $f \mapsto (Kf)(p)$ with respect to
$\norm{\cdot}_{\lp}$ is:
\[
  b_p
  = \sup_{\norm{f} = 1} \abs{(Kf)(p)}
  = \sup_{\norm{f} = 1} \abs[\Big]{\sum_q K(p,q) f(q)}
  \leq \left(\sum_q \abs{K(p,q)}^2\right)^{1/2}
  = p^{-\alpha}\, \norm{h}_{L^2},
\]
where the last equality follows from:
\[
  \sum_{q \in \PP} \abs{K(p,q)}^2
  = p^{-2\alpha} \sum_q q^{-2\alpha} \abs{h(\log p - \log q)}^2
  \lesssim p^{-2\alpha}\, \norm{h}_{L^2}^2. \qedhere
\]
\end{proof}

\begin{appremark}[Implementation binding]
\cref{prop:alpha-bp} is the precise statement of the
\texttt{ZMOperatorStack.\_b\_map} computation documented in
\texttt{pirtm/zm/operator\_stack.py}:
\[
  \texttt{self.\_b\_map[p]}
  \;\longleftrightarrow\;
  p^{-\texttt{self.alpha}} \cdot \norm{h}_{L^2}.
\]
This is the co-dependency bridge: \texttt{certify\_state(b\_p, ...)} receives
the output of \texttt{build\_prime\_block()}, which computes
\texttt{\_b\_map} using this formula.
\end{appremark}

% ============================================================
\section{Spin Foam Identification: Area Eigenvalue as Prime Weight}
\label{app:spinfoam}
% ============================================================

\begin{appproposition}[Spin foam as prime-gated instantiation]
\label{prop:spinfoam}
Let $\hat{A}$ be the LQG area operator with eigenvalues
$a_j = \ell_P^2 \sqrt{j(j+1)}$ for $j \in \tfrac{1}{2}\ZZ_{\geq 0}$.
Define the identification map $\iota \colon \tfrac{1}{2}\ZZ_{\geq 0} \to \PP$
by enumerating spin values in order and mapping to the $n$-th prime:
$\iota(j_n) = p_n$. Then under the substitution
\[
  p^{-\sigma}
  \;\longleftrightarrow\;
  \ell_P^{-2}\, a_{\iota^{-1}(p)}
  \;=\;
  \sqrt{j(j+1)}\big|_{j = \iota^{-1}(p)},
\]
the prime block $A = D_\sigma + K$ acts on the spin network Hilbert space
$\mathcal{H}_{\mathrm{spin}} \cong \lp$ (under $\iota$), and the ZM
certification protocol applies without modification to spin foam amplitudes.
\end{appproposition}

\begin{proof}
The identification $\mathcal{H}_{\mathrm{spin}} \cong \lp$ follows from
the separability of both spaces and the countability of the spin spectrum
$\{j_n\}_{n \geq 0}$. The bijection $\iota$ is order-preserving by
construction. The diagonal of $D_\sigma$ in spin foam coordinates becomes
the area spectrum weighting $a_j / \ell_P^2$, with $\sigma$ playing the role
of the area-to-weight exponent. The HS off-diagonal kernel $K$ encodes
spin foam vertex amplitudes coupling adjacent spin labels.
\end{proof}

% ============================================================
\section{Consolidated Norm Bound Table}
\label{app:bounds-table}
% ============================================================

The following table collects all operator norm bounds derived in this appendix
for the reference parameter set
$\sigma = 0.8$, $\alpha = 0.8$, $\norm{h}_{L^2} = 1$,
$\norm{m}_{L^\infty} = 0.5$, $\norm{E}_{\op} = 0.3$.

\begin{longtable}{llll}
\toprule
\textbf{Operator} & \textbf{Bound Formula} & \textbf{Reference}
  & \textbf{Numerical Value} \\
\midrule
$\norm{D_\sigma}_{\op}$ & $2^{-\sigma}$ & \cref{thm:A-norm}
  & $0.574$ \\[2pt]
$\norm{K}_{\HS}$ & $\zeta_\PP(2\alpha)\norm{h}_{L^2}$ & \cref{thm:hs}
  & $0.452$ \\[2pt]
$\norm{K}_{\op}$ & $\leq \norm{K}_{\HS}$ & \cref{thm:A-norm}
  & $\leq 0.452$ \\[2pt]
$\norm{A}_{\op}$ & $2^{-\sigma} + \zeta_\PP(2\alpha)\norm{h}_{L^2}$
  & \cref{thm:A-norm} & $\leq 1.026$ \\[2pt]
$\sgap^*(D_\sigma)$ & $2^{-\sigma} - 3^{-\sigma}$
  & \cref{lem:sgap-lower} & $0.574 - 0.431 = 0.143$ \\[2pt]
$\norm{C}_{\op}$ & $\norm{m}_{L^\infty}$ & \cref{thm:U-norm} & $0.500$ \\[2pt]
$\norm{U}_{\op}$ & $\norm{A}_{\op}+\norm{C}_{\op}+\norm{E}_{\op}$
  & \cref{thm:U-norm} & $\leq 1.826$ \\[2pt]
$\gamma_{\mathrm{eff}}(t)$ & $\leq \gamma(1-\alpha_0(t))$
  & \cref{thm:xi-contraction} & $<\gamma$ as $t\to\infty$ \\[2pt]
$\GapLB(w)$ & $\inf_\theta\delta_S(\theta) - 2\sum_p\abs{w_p}b_p$
  & \cref{thm:gaplb-full} & (parameter-dependent) \\[2pt]
$\SlopeUB(w)$ & $L\cdot w_{\mathrm{budget}}$
  & \cref{thm:gaplb-full} & (parameter-dependent) \\[2pt]
$b_p$ & $p^{-\alpha}\norm{h}_{L^2}$
  & \cref{prop:alpha-bp} & $p^{-0.8}$ \\
\bottomrule
\caption{Consolidated operator norm bounds (reference parameters $\sigma=0.8$,
$\alpha=0.8$, $\norm{h}_{L^2}=1$).}
\end{longtable}

\nocite{*}
\bibliographystyle{unsrt}  
\bibliography{references}
\end{document}